\chapter[Sermon 29. Here Is Described the Commendation and Hymn Said unto Christ of the Angels and All Creatures. Etc.]{Here Is Described the Commendation and Hymn Said unto Christ of the Angels and All Creatures. Etc.}
\label{sermon:029}
\markright{Sermon 29. Here Is Described the Commendation and Hymn Said unto ...}

And I looked, and I heard the voice of countless angels around the throne, around the beast and the elders, and I heard thousands of thousands saying with a mighty voice: “Worthy is the Lamb who was slain to receive power, and riches, and wisdom, and strength, and honor, and glory, and blessing.” And every creature that is in heaven, and on the earth, and under the earth, and in the sea, and all that are in them, heard. Saying: “Blessing, honor, glory, and power be to him who sits on the throne, and to the Lamb forever and ever.” And the four living creatures said, “Amen.” And the twenty-four elders fell on their faces and worshiped him who lives forever and ever.

Fourthly, the angels of God now come to the elders and the beasts, that is, to God’s most excellent creatures, and together with them offer praise to God and the Lamb. This is surely an example for us, as I often say and repeat, so that we may understand what is fitting for us as well.

\index{Daniel (Book of)}\index{Hebrews, Epistle to the}Of Angels David in Psalm 4, speaking among other things: which says he makes his Angels spirits, and his ministers a flame of fire. He testifies therefore that the Angels were made or created by God. By their substance he calls them spirits, and by a parable a flame of fire, which is pure, bright, most swift, piercing, and burning. Therefore, in their sort and manner, the Angelical spirits are altogether such: whom by their office he calls ministers, to wit, of God and man. For Paul also to the Hebrews, bringing this same place of David, asks, “Are they not all ministering spirits, sent forth to serve those who will inherit salvation?” truly understanding men. These things teach us to judge rightly of Angels, and that no man should worship ministers, or any, however excellent creatures they may be, for their godly gifts. Neither, in deed, can the Angels or Saints abide being worshipped. Here, without doubt, they attribute all glory to God and to the Lamb, to God three and one, that all we should do the like. Here is also declared the place wherein the Angels were, about the throne, about the beasts, and about the Elders. Therefore, they guarded all these places round about as it were a guard. Daniel in times past saw things not much unlike these in chapter 7. Certainly, they stand like ministers and servants, ready to do service. Angels are said also to be about the godly upon earth, and to attend upon the salvation and ministry of men. In Psalm 34, David sings, that being afflicted he called upon the Lord, and the Lord heard him, and from all his troubles he delivered him. The Angel of the Lord pitches his tents about those who fear him, and he delivers them. And not much unlike things you may read in Psalm 91. And you shall note here, that those who are afflicted do call upon the Lord, and not the Angels: and that the Lord does hear, and deliver, and for the working thereof uses the ministry of Angels, as his ministers. And just as no man of sound mind reveres, calls upon, and worships the son, because God by him gives great benefits to men, so no man honors, calls upon, and worships Angels, because God uses their ministry in delivering men.

Now he also states the number of angels, but a certain number for an uncertain one, thousands of thousands for innumerable. He alludes in the meantime to that saying of Daniel in chapter 7: “Ten thousands ministered to him, and ten thousand times ten assisted him.” We are accustomed to consider the power of kings as a great and innumerable army. What, then, shall we think of the power of our God, who is the God of hosts, and whom not only innumerable legions of angels, but all creatures serve? And what an excellent praise is sung together by so many blessed spirits?

After this, the proper duty of the angels is addressed: they sing praises to God and commend God’s work, and that with a loud voice. It matters little whether you sing praises to the Lord with a low or a high voice; but because those who cry with a loud voice are often deeply moved, whether overwhelmed with great sorrow or rejoicing with great gladness, therefore we shall praise God with a loud voice, if with a steadfast spirit and with the inward affection of the heart we praise God.

The evangelical hymn is now included, which agrees in all things with the Hymns of the Beasts and Elders. For they celebrate the Lamb, that is to say the Son, which as he is the savior alone, so hath he deserved to receive all power and glory, and to govern all things: as is said before.

\index{Primasius}And seven things the Angles attribute to the label, that is to Jesus Christ our Lord, on the right hand of the Father. First, power, namely godly, almighty, vivificative, and conservative. Of this I spoke also before. Secondly, riches. For he is rich, as the Apostle says, for all who call upon him. Romans 10. And Primasius says: Christ himself is the treasure of all good things. For \textit{schaddais} is the sufficiency of all goods of the mind and body. And if it be lawful to attribute a profane word unto God, he is very Saturn, fulfilling all creatures. And the Angles commend Christ so, who would think that men should strive so to acquire things for themselves, as though they themselves could fill their own desires? Then they attribute to Christ wisdom, namely godly and great. For the Son is the wisdom of the Father, a topic much treated by Solomon. By this wisdom he can rule all things by the most agreeable and best government. Who shall say? Thus it should have been done. The wisdom of God has most beautifully and well made all things from the beginning, so that our reason can justly blame nothing: what thing shall we blame now in the universal government of Christ? They ascribe unto Christ also strength to execute truly such things as he has most wisely ordained, finally strength to defend his, and to subdue the adversaries. For he is almighty. Such things as follow, honor, glory, and thanksgiving, are declared before, what they are, and of what force; saving that the thing he called first [illegible] he calls now [illegible], blessing, praise, and giving of thanks.

\index{Augustine, Saint}This hymn, sung in praise of Christ, teaches that Christ is truly God, of the same substance and coequal with the Father, greater than angels, indeed the Lord of angels, whom the angels themselves also worship, as Paul declared in the first letter to the Hebrews. Therefore, those who place angels above Christ are refuted. The heretics called Angelicals—that is, worshippers of angels—are refuted. The angels themselves here rebuke their error: rightly are they considered heretics by Augustine. If riches, glory, and honor are due to Christ alone, and he surpasses them, why are these communicated to creatures? Otherwise, we admonish all faithful people to esteem angels highly, to acknowledge and marvel at the benefits of God working through them, and to love them as brethren, companions, and coheirs of the same salvation; much less to despise or blame them. Of this I will speak more at another time.

Hitherto he has recited the excellent praises, the celebratory verses or hymns, of God’s excellent creatures, especially of the Elders and generally of all, moreover of Angels also, spoken to Christ our redeemer and prince. And yet not content with these, he adds moreover in the fifth place, the agreement, and praise, and submission of all the creatures in the world, so that if we are not moved by the excellent example of the excellent creatures, the Elders and Angels, at the last we might be ashamed, seeing all creatures of their own accord do their duty. For man sins; he is Lord of all, and all things were created for him. How, I pray you, shall he offend more heinously against God, who has made him Lord over all, than if by his hardness, ingratitude, and malice he not only does not do his own duty, but is rather inferior to all creatures; as he who alone contends with God, and does not attribute to him due praise? Therefore this example greatly excites man that he should submit himself unto God, and give God the whole glory; and in no wise strive with God, nor complain of anything. But mark, I pray you, with how diligent a division of things he encompasses all creatures, excluding none, the Devil only excepted, when he recounts the creatures that are in Heaven, that are in Earth, that are under the Earth, and in the Sea; finally, he adds, and all that are therein. Therefore, if all created things do celebrate and worship him that sits in the Throne and the Lamb, and submit themselves unto him, is it not a shame, yea and a soul shame, that man alone, Lord of all, should revolt to the sworn enemy of God, the Devil, and with him to expostulate with God, to teach and blame and find fault with his judgments and governments, to complain of his works and will?

You marvel, I well know, how all creatures, even those without reason and feeling, can praise God. However, this figure of personification—the attributing of human qualities to non-human things—is very common among the prophets, and especially with David, who says, “Praise him, sun and moon; praise him, bright stars.” “Praise the Lord from the earth, sea monsters and all deep places; fire and hail, snow and ice.” And by such ways of speaking, the prophets would encourage and stir up people to praise God, saying that creatures which have no life, do after their manner praise God; see that you after your manner do praise God in hymns and spiritual psalms. Indeed, David shows a plain reason why he commands bodies that have no life to praise God: “Let them praise the name of the Lord,” he says, “because he commanded, and they were created.” As though he should say: they are his creatures, and in that they remain, they have it of him; therefore let them make the name of God glorious, as of their maker and preserver. And he signified also the manner of praising, where he adds, “He hath ordained them, that they should endure forever; he hath appointed a decree, and it has not passed away.” As if he should have said: Where they neglect no part of those things for which they were made, but are ready in their place, order, and time, and do their duty exceedingly well, do they not preach unto men the wonderful wisdom and power of God? For in another psalm also, David says, “The heavens show forth the glory of God, and the firmament declares the works of his hands.” Thus I say, the creatures without life do praise and commend the name of God to men, when they are moved, work wonderfully, and obediently do the thing to which they are appointed.

The Hymn of all creatures, similar to that which was of the beasts, Elders, and Angels, is here also suitably described, though briefly. But where it contains nothing that has not already been declared, I will not weary the gentle listeners by repeatedly speaking the same things. However, one thing seems chiefly to be observed: that they join him who sits on the throne and the Lamb together, thus acknowledging the Son to be coequal with the Father, and that both are to be worshipped with like honor, and with like praises celebrated and commended. They attribute uniquely to the Lamb dominion or kingdom, for he received the book of the Father, as has been declared before: namely, all power and authority to govern all things.

The four beasts sing to it Amen, either confirming the hymn of the creatures, or declaring their consent with them. This is so that we may, with one mind, pray together and praise God, blessed forevermore. With these are moreover refuted the dispensations of men. The Lord allows the concord and agreement of men, and requires it utterly, especially in prayers and godly praises. For he commands in the Gospel to lay down your offering, which you would offer, if you remember any discord between you and your brother, to go to him, and to renew friendship, and then to return to your offering: which in the prophets is called an abomination, if it be offered by minds possessed with rancor and malice. \&c.

Finally, the elders fall down again and worship him who lives forever, no doubt so that all of us on earth might be moved to obedience. For if these things are done in heaven by the blessed spirits, what, I pray you, is fitting for us to do here on earth? And note that they are said to worship him who lives forever, who nevertheless also fell down before the Lamb and before the Throne, from which the Spirit proceeded, and where he who sits is seated; from which we gather that the Father, the Son, and the Holy Spirit are indeed distinct in persons, yet these three are not three Gods, but one God living forever.

And truly, this notable vision and treatise may serve as an effective remedy against various poisons of heresies, especially those of the Arians and Servetans, or rather, [?perdetans], moreover against diverse and curious disputations and temptations concerning the works, judgments, and providence of God. If we are wise, we will obediently submit ourselves to the living God with all the creatures and saints of God, worshipping him, and with the prophet crying: “You are a just Lord in all your ways, and holy in all your works. You have created us; all things are yours. You govern all things in the best order. You love humankind. You have given us your Son. You, by your Son our redeemer, govern all things uprightly. We worship you, the Father, the Son, and the Holy Spirit, one very God. To you is due the kingdom, honor, and glory forever and ever. Amen.”
